\subsection{同伦箭的缩并}

设 H 为一个箭图，G 为一个群胚，\(\varphi\) 与 \(\psi\) 为从 H 到 G 的两个箭图态射，并设 \(h\colon \mathrm{Som}(\mathrm{H})\to \mathrm{Fl}(\mathrm{G})\) 为连接 \(\varphi\) 与 \(\psi\) 的一个同伦。设 G′ 为由 G 通过缩并 \(h\) 的像中的箭而得到的群胚（II, 第 175 页, 第 11 号），并设 \(\beta\colon\mathrm{G}\to\mathrm{G}'\) 为典范态射。

以 \(\Gamma\) 表示箭图 (Som(G), Som(H), Som(φ), Som(ψ))。根据同伦的定义，偶 (Id \(_{\text{Som(G)}}\) , \(h\) ) 是从 \(\Gamma\) 到 G 的一个箭图态射；它唯一地扩张为群胚态射 \(\eta\) : Grp(Γ) → G。根据构造，G′ 的顶点集是箭图 \(\Gamma\) 的连通分支所成的集。

在本节（\(n^{o}\)）的余下部分，我们将假设群胚 G 是传递的。由 II, 第 170 页的注 1，这等价于假设群胚 G′ 是传递的。我们还固定 G 的一个顶点 \(a_{0}\) 。

以 \(\widetilde{\Gamma}\) 表示与箭图 \(\Gamma\) 相伴的图（参看 II, 第 156 页）。长度 \(\geqslant 1\) 的回路（\(\widetilde{\Gamma}\) 中的）所成的集与集 \(\Omega(\widetilde{\Gamma})\) （有限序列 \((z_1, \ldots, z_n)\) 所成之集）等同，其中 \(n\) 是满足 \(n \geqslant 1\) 的整数，而 \(z_1, \ldots, z_n\) 是 \(\mathrm{Fl}(\widetilde{\Gamma})\) 的元素，满足 \(t(z_i) = o(z_{i+1})\) （对每个整数 \(i\) ，满足 \(1 \leqslant i < n\) ）以及 \(t(z_n) = o(z_1)\) 。

设 \(\mathbf{z} = ((b_{1}, \varepsilon_{1}), \ldots, (b_{n}, \varepsilon_{n}))\) 为 \(\Omega(\widetilde{\Gamma})\) 的一个元素。元素

\[
\operatorname{Int} (g) (\eta (\mathbf {z})) = g h (b _ {1}) ^ {\varepsilon_ {1}} \dots h (b _ {n}) ^ {\varepsilon_ {n}} g ^ {- 1}
\]

在 \(G_{a_{0}}\) 中的共轭类不依赖于 G 的、连接顶点 \(a_{0}\) 与 \((b_{1},\varepsilon)\) 的起点的箭 g 的选取。我们以 \(c(\mathbf{z})\) 表示此共轭类。

命题 1. —— 群同态 \(\beta_{a_0} \colon \mathrm{G}_{a_0} \to \mathrm{G}'_{\beta(a_0)}\) 是满射的，其核是 \(\mathrm{G}_{a_0}\) 的、包含共轭类 \(c(\mathbf{z})\) （对 \(\mathbf{z} \in \Omega(\widetilde{\Gamma})\) ）的最小子群。

设 K 为 G 的最小的正规子群胚，其箭集包含 h 的像。态射 \(\mathrm{G}_{a_{0}} \rightarrow \mathrm{G}^{\prime}_{\beta(a_{0})}\) 是满射的，且其核等于 \(K_{a_{0}}\) （II, 第 170 页, 命题 2）。于是命题由 II, 第 176 页的命题 8 得出。

定义 1. —— 称 \(\Omega(\widetilde{\Gamma})\) 的一个子集 Z 为正规的（相对于偶 \((\varphi, \psi)\) ），如果它满足下列性质：

\begin{enumerate}
\def\labelenumi{(\roman{enumi})}
\item
  对 Γ 的每条箭 z，有 \((z,\overline{z})\in\mathbb{Z}\) ；
\item
  对每个 \((z_{1},\ldots ,z_{n})\in \mathbb{Z}\) ，我们有 \((\overline{z}_n,\dots ,\overline{z}_2,\overline{z}_1)\in \mathbb{Z}\) 以及 \((z_{n},z_{1},\ldots ,z_{n - 1})\in \mathbb{Z}\) ；
\item
  设 \(\mathbf{z} = (z_1, \ldots, z_n)\) 与 \(\mathbf{z}' = (z_1', \ldots, z_m')\) 为 \(Z\) 中满足 \(t(z_n) = o(z_1')\) 的元素。令 \(\mathbf{z}\mathbf{z}' = (z_1, \ldots, z_n, z_1', \ldots, z_m')\) 。若 z、z′、zz′ 三者中有两个属于 Z，则第三个也属于 Z；
\item
  对 H 的每条箭 \(f\) ，令 \(\widetilde{\varphi}(f,1) = \widetilde{\psi}(f,-1) = \varphi(f)\) 以及 \(\widetilde{\varphi}(f,-1) = \widetilde{\psi}(f,1) = \psi(f)\) 。设 \(n\) 为一个整数 \(\geqslant 1\) ，并设 \((f_1,\varepsilon_1),\ldots,(f_n,\varepsilon_n)\) 为 \(\mathrm{Fl}(\mathrm{H})\times \{-1,1\}\) 的元素的一个序列，满足 \(\widetilde{\psi}(f_i,\varepsilon_i) = \widetilde{\varphi}(f_{i+1},\varepsilon_{i+1})\) （对 \(1\leqslant i < n\) ）以及 \(\widetilde{\psi}(f_n,\varepsilon_n) = \widetilde{\varphi}(f_1,\varepsilon_1)\) ；以 \(a_i\) 表示 \(f_i\) 的起点，以 \(b_i\) 表示其靶。\(((a_1,\varepsilon_1),\ldots,(a_n,\varepsilon_n))\) 属于 Z 当且仅当 \(((b_1,\varepsilon_1),\ldots,(b_n,\varepsilon_n))\) 属于 Z。
\end{enumerate}

\(\Omega(\widetilde{\Gamma})\) 中任意一族正规子集（相对于 \((\varphi, \psi)\) ）的交仍是正规的。特别地，对给定子集 \(Z\) （\(\Omega(\widetilde{\Gamma})\) 的），存在包含它的最小正规子集。

命题 2. —— 若 \(N\) 是 \(G_{a}\) 的一个正规子群，则元素 \(\mathbf{z} \in \Omega(\widetilde{\Gamma})\) （满足 \(c(\mathbf{z})\) 包含于 \(N\) ）所成的集是 \(\Omega(\widetilde{\Gamma})\) 的一个正规子集。

以 Z 表示元素 \(\mathbf{z} \in \Omega(\widetilde{\Gamma})\) （满足 \(c(\mathbf{z}) \subset \mathrm{N}\) ）所成的集。

设 \(z \in \mathrm{Fl}(\widetilde{\Gamma})\) 。根据定义有 \(c(z, \overline{z}) = \{e_a\}\) ，从而 \((z, \overline{z}) \in \mathbf{Z}\) 。设 \((z_1, \ldots, z_n) \in \mathbf{Z}\) 。根据 \(c\) 的定义，共轭类 \(c(z_1, \ldots, z_n)\) 等于 \(c(z_n, z_1, \ldots, z_{n-1})\) ，并且由 \(c(\overline{z}_n, \ldots, \overline{z}_1)\) 的元素的逆所组成。这表明 \(\mathbf{Z}\) 满足条件 (ii)。

采用条件 (iii) 的记号，可以选取元素 \(u \in c(\mathbf{z})\) 、\(v \in c(\mathbf{z}')\) 与 \(w \in c(\mathbf{z}\mathbf{z}')\) ，使得 uv = w。若元素 u、v 与 w 三者中有两个属于 N，则第三个亦然，因此 N 满足 (iii)。

采用 (iv) 中的记号，对每个满足 \(1 \leqslant i \leqslant n\) 的整数 i，有关系式

\[
  \varphi (f _ {i}) h (b _ {i}) = h (a _ {i}) \psi (f _ {i}),
\]

因为 h 是连接 \(\varphi\) 与 \(\psi\) 的一个同伦。此等式也可写成

\[
\widetilde {\varphi} (f _ {i}, \varepsilon_ {i}) h (b _ {i}) ^ {\varepsilon_ {i}} = h (a _ {i}) ^ {\varepsilon_ {i}} \widetilde {\psi} (f _ {i}, \varepsilon_ {i}).
\]

考虑到关系式 \(\widetilde{\psi}(f_i,\varepsilon_i) = \widetilde{\varphi}(f_{i + 1},\varepsilon_{i + 1})\) （对 \(1\leqslant i <   n\) ）以及 \(\widetilde{\psi} (f_n,\varepsilon_n) = \widetilde{\varphi} (f_1,\varepsilon_1)\) ，我们推出

\[
\widetilde {\varphi} (f _ {1}, \varepsilon_ {1}) h (b _ {1}) ^ {\varepsilon_ {1}} \dots h (b _ {n}) ^ {\varepsilon_ {n}} = h (a _ {1}) ^ {\varepsilon_ {1}} \dots h (a _ {n}) ^ {\varepsilon_ {n}} \widetilde {\varphi} (f _ {1}, \varepsilon_ {1}),
\]

从而共轭类 \(c((a_{1},\varepsilon_{1}),\ldots,(a_{n},\varepsilon_{n}))\) 与 \(c((b_{1},\varepsilon_{1}),\ldots,(b_{n},\varepsilon_{n}))\) 相等。这表明 Z 满足条件 (iv)，命题证明完毕。

推论. —— 设 Z 为 \(\Omega(\widetilde{\Gamma})\) 的一个子集。为使共轭类 \(c(\mathbf{z})\) （对 \(\mathbf{z} \in \mathbf{Z}\)）生成典范同态 \(\beta_{a_0}: \mathrm{G}_{a_0} \to \mathrm{G}'_{\beta(a_0)}\) 的核，只需 \(\Omega(\widetilde{\Gamma})\) 中包含 Z 的最小正规子集等于 \(\Omega(\widetilde{\Gamma})\)。

设 \(N\) 为 \(G_{a_{0}}\) 中包含 \(c(\mathbf{z})\) （对所有 \(z \in Z\) ）的最小子群；它是正规的。设 \(Z'\) 为元素 \(z\) （\(\Omega(\widetilde{\Gamma})\) 中的，满足 \(c(\mathbf{z}) \in \mathrm{N}\) ）所成的集。由命题 2，\(Z'\) 是 \(\Omega(\widetilde{\Gamma})\) 的一个正规子集。它包含 Z。根据假设，它于是等于 \(\Omega(\widetilde{\Gamma})\) 。于是由 II, 第 197 页的命题 1 得出，N 是同态 \(\beta_{a}\) 的核。

